% ============================================================================
% CHAPTER 9: REINFORCEMENT LEARNING FRAMEWORK
% ============================================================================

\chapter{Reinforcement Learning Framework}

\section{RL System Overview}

\subsection{Introduction to RL in Classification}

The Reinforcement Learning (RL) framework optimizes the pharmacovigilance classification system by learning from real-world data and feedback to dynamically adjust classification strategies.

\textbf{Key Objectives:}
\begin{itemize}
    \item \textbf{Maximize Accuracy}: Achieve 92%+ classification accuracy
    \item \textbf{Minimize Cost}: Reduce LLM API costs by 65%
    \item \textbf{Optimize Speed}: Reduce average processing time by 44%
    \item \textbf{Maintain Stability}: Ensure consistent performance across domains
    \item \textbf{Domain Specialization}: Adapt to specific medical domains
\end{itemize}

\subsection{RL vs. Traditional ML}

\begin{table}[h]
\centering
\caption{RL vs. Traditional ML Comparison}
\begin{tabularx}{\textwidth}{lXX}
\toprule
\textbf{Aspect} & \textbf{Traditional ML} & \textbf{Reinforcement Learning} \\
\midrule
Training Data & Fixed dataset & Continuous interaction \\
Feedback & Immediate labels & Delayed rewards \\
Adaptation & Static model & Dynamic adaptation \\
Optimization & Minimize loss function & Maximize long-term reward \\
Decision Making & Single prediction & Sequential decisions \\
Exploration & N/A & Exploration vs. Exploitation \\
Use Case & Classify documents & Optimize classification strategy \\
\bottomrule
\end{tabularx}
\end{table}

\section{Multi-Objective RL Architecture}

\subsection{State Space (47 Dimensions)}

The RL agent observes a 47-dimensional state vector for each document.

\textbf{State Components:}
\begin{lstlisting}[style=pythonstyle]
# rl_optimization/rl_training/rl_agent.py
class DocumentFeatureExtractor:
    def extract_state(self, document: dict) -> np.ndarray:
        """Extract 47-dimensional state vector"""
        
        # 1. Document Features (15 dimensions)
        doc_features = [
            len(document.get('content', '')),  # Document length
            len(document.get('abstract', '')),  # Abstract length
            document.get('page_count', 0),  # Page count
            len(document.get('mesh_terms', [])),  # MeSH term count
            len(document.get('keywords', [])),  # Keyword count
            self._calculate_medical_term_density(document),
            self._calculate_patient_indicator_score(document),
            self._calculate_drug_mention_score(document),
            self._calculate_adverse_event_score(document),
            self._calculate_temporal_marker_score(document),
            self._has_structured_sections(document),  # 0 or 1
            self._publication_year(document),
            self._journal_impact_factor(document),
            self._author_count(document),
            self._citation_count(document)
        ]
        
        # 2. Classification History (10 dimensions)
        history_features = [
            self._get_recent_accuracy(window=100),
            self._get_recent_cost(window=100),
            self._get_recent_speed(window=100),
            self._get_recent_confidence_variance(),
            self._get_classification_distribution('ICSR'),
            self._get_classification_distribution('Safety_Information'),
            self._get_classification_distribution('Non-ICSR'),
            self._get_classification_distribution('Rejected'),
            self._get_cascade_activation_rate(),
            self._get_llm_utilization_rate()
        ]
        
        # 3. System Metrics (12 dimensions)
        system_features = [
            self._get_queue_depth(),
            self._get_cpu_usage(),
            self._get_memory_usage(),
            self._get_request_rate(),
            self._get_error_rate(),
            self._get_api_latency(),
            self._get_database_load(),
            self._get_concurrent_users(),
            self._get_time_of_day(),  # 0-23
            self._get_day_of_week(),  # 0-6
            self._get_system_uptime(),
            self._get_service_health_score()
        ]
        
        # 4. Domain Indicators (10 dimensions)
        domain_features = [
            self._classify_domain(document),  # One-hot: cardio, neuro, etc.
            self._get_domain_confidence(),
            self._get_domain_complexity(),
            self._get_regulatory_relevance(),
            self._is_clinical_trial(document),
            self._is_case_report(document),
            self._is_meta_analysis(document),
            self._has_statistical_data(document),
            self._language_indicator(document),
            self._source_reliability_score(document)
        ]
        
        # Combine all features
        state = np.concatenate([
            doc_features,
            history_features,
            system_features,
            domain_features
        ])
        
        # Normalize to [0, 1] range
        return self._normalize_state(state)
\end{lstlisting}

\subsection{Action Space (12 Actions)}

The RL agent can take 12 discrete actions to modify classification behavior.

\begin{table}[h]
\centering
\caption{RL Action Space}
\begin{tabularx}{\textwidth}{lXl}
\toprule
\textbf{Action ID} & \textbf{Description} & \textbf{Effect} \\
\midrule
0 & Increase rule\_based weight by 0.1 & More reliance on rules \\
1 & Decrease rule\_based weight by 0.1 & Less reliance on rules \\
2 & Increase BERT weight by 0.1 & More reliance on BERT \\
3 & Decrease BERT weight by 0.1 & Less reliance on BERT \\
4 & Increase LLM weight by 0.1 & More reliance on LLM \\
5 & Decrease LLM weight by 0.1 & Less reliance on LLM \\
6 & Increase rule confidence threshold by 0.05 & Stricter rule requirement \\
7 & Decrease rule confidence threshold by 0.05 & Lenient rule requirement \\
8 & Increase BERT confidence threshold by 0.05 & Stricter BERT requirement \\
9 & Decrease BERT confidence threshold by 0.05 & Lenient BERT requirement \\
10 & Enable cascade (skip lower models if high confidence) & Speed optimization \\
11 & Disable cascade (always run all models) & Accuracy optimization \\
\bottomrule
\end{tabularx}
\end{table}

\textbf{Action Implementation:}
\begin{lstlisting}[style=pythonstyle]
class MultiObjectiveRLAgent:
    def __init__(self):
        self.state_dim = 47
        self.action_dim = 12
        
        # Current configuration (modified by actions)
        self.config = {
            'rule_weight': 0.30,
            'bert_weight': 0.48,
            'llm_weight': 0.22,
            'rule_threshold': 0.85,
            'bert_threshold': 0.75,
            'cascade_enabled': True
        }
    
    def apply_action(self, action: int):
        """Apply action to modify classification configuration"""
        if action == 0:
            self.config['rule_weight'] = min(1.0, self.config['rule_weight'] + 0.1)
        elif action == 1:
            self.config['rule_weight'] = max(0.0, self.config['rule_weight'] - 0.1)
        # ... (similar for other actions)
        
        # Normalize weights to sum to 1.0
        total = (self.config['rule_weight'] + 
                 self.config['bert_weight'] + 
                 self.config['llm_weight'])
        self.config['rule_weight'] /= total
        self.config['bert_weight'] /= total
        self.config['llm_weight'] /= total
\end{lstlisting}

\subsection{Reward Function (Multi-Objective)}

The RL agent optimizes for 5 competing objectives with dynamic weights.

\textbf{Reward Calculation:}
\begin{lstlisting}[style=pythonstyle]
def calculate_reward(self, classification_result: dict, ground_truth: str) -> float:
    """Calculate multi-objective reward"""
    
    # 1. Accuracy Reward (weight: 0.40)
    accuracy_reward = 1.0 if classification_result['class'] == ground_truth else -0.5
    
    # 2. Cost Efficiency Reward (weight: 0.25)
    # Reward for avoiding expensive LLM calls
    cost = classification_result.get('cost', 0)
    max_cost = 0.001  # Max expected cost
    cost_reward = 1.0 - (cost / max_cost)
    
    # 3. Speed Reward (weight: 0.15)
    # Reward for fast processing
    processing_time = classification_result.get('processing_time_ms', 0)
    max_time = 5000  # 5 seconds
    speed_reward = 1.0 - (processing_time / max_time)
    
    # 4. Stability Reward (weight: 0.10)
    # Reward for consistent confidence scores
    confidence = classification_result.get('confidence', 0)
    confidence_variance = self._get_recent_confidence_variance()
    stability_reward = 1.0 - confidence_variance
    
    # 5. Domain Specialization Reward (weight: 0.10)
    # Reward for correct classification in specialized domains
    domain = self._classify_domain(classification_result['document'])
    domain_difficulty = self._get_domain_difficulty(domain)
    if classification_result['class'] == ground_truth:
        domain_reward = domain_difficulty  # Higher reward for difficult domains
    else:
        domain_reward = -domain_difficulty
    
    # Weighted sum
    total_reward = (
        0.40 * accuracy_reward +
        0.25 * cost_reward +
        0.15 * speed_reward +
        0.10 * stability_reward +
        0.10 * domain_reward
    )
    
    # Clip to [-1, 1] range
    return np.clip(total_reward, -1.0, 1.0)
\end{lstlisting}

\begin{figure}[h]
\centering
\fbox{\textit{[SCREENSHOT PLACEHOLDER: RL Reward Components - Pie Chart]}}
\caption{Multi-Objective Reward Breakdown}
\label{fig:rl-reward-breakdown}
\end{figure}

\section{RL Algorithms Implementation}

\subsection{Actor-Critic Architecture}

The RL agent uses a dual-network Actor-Critic architecture.

\textbf{Neural Network Architecture:}
\begin{lstlisting}[style=pythonstyle]
import torch
import torch.nn as nn

class ActorNetwork(nn.Module):
    """Policy network: state -> action probabilities"""
    def __init__(self, state_dim=47, action_dim=12, hidden_dim=256):
        super().__init__()
        self.network = nn.Sequential(
            nn.Linear(state_dim, hidden_dim),
            nn.ReLU(),
            nn.Dropout(0.2),
            nn.Linear(hidden_dim, hidden_dim),
            nn.ReLU(),
            nn.Dropout(0.2),
            nn.Linear(hidden_dim, action_dim),
            nn.Softmax(dim=-1)  # Action probabilities
        )
    
    def forward(self, state):
        return self.network(state)

class CriticNetwork(nn.Module):
    """Value network: state -> expected return"""
    def __init__(self, state_dim=47, hidden_dim=256):
        super().__init__()
        self.network = nn.Sequential(
            nn.Linear(state_dim, hidden_dim),
            nn.ReLU(),
            nn.Dropout(0.2),
            nn.Linear(hidden_dim, hidden_dim),
            nn.ReLU(),
            nn.Dropout(0.2),
            nn.Linear(hidden_dim, 1)  # Single value output
        )
    
    def forward(self, state):
        return self.network(state)
\end{lstlisting}

\subsection{Training Algorithm (A2C)}

Advantage Actor-Critic (A2C) for stable policy gradient updates.

\begin{lstlisting}[style=pythonstyle]
class A2CAgent:
    def __init__(self, state_dim=47, action_dim=12):
        self.actor = ActorNetwork(state_dim, action_dim)
        self.critic = CriticNetwork(state_dim)
        
        self.actor_optimizer = torch.optim.Adam(self.actor.parameters(), lr=0.001)
        self.critic_optimizer = torch.optim.Adam(self.critic.parameters(), lr=0.001)
        
        self.gamma = 0.99  # Discount factor
        self.entropy_coef = 0.01  # Encourage exploration
    
    def select_action(self, state: np.ndarray, training=True):
        """Select action using policy network"""
        state_tensor = torch.FloatTensor(state).unsqueeze(0)
        action_probs = self.actor(state_tensor)
        
        if training:
            # Sample from distribution (exploration)
            dist = torch.distributions.Categorical(action_probs)
            action = dist.sample()
            log_prob = dist.log_prob(action)
            entropy = dist.entropy()
            return action.item(), log_prob, entropy
        else:
            # Greedy selection (exploitation)
            return torch.argmax(action_probs).item(), None, None
    
    def train_step(self, trajectory: list):
        """Train on single trajectory (episode)"""
        states, actions, rewards, next_states, dones = zip(*trajectory)
        
        # Convert to tensors
        states = torch.FloatTensor(states)
        actions = torch.LongTensor(actions)
        rewards = torch.FloatTensor(rewards)
        next_states = torch.FloatTensor(next_states)
        dones = torch.FloatTensor(dones)
        
        # Compute values
        values = self.critic(states).squeeze()
        next_values = self.critic(next_states).squeeze()
        
        # Compute TD targets and advantages
        td_targets = rewards + self.gamma * next_values * (1 - dones)
        advantages = td_targets - values
        
        # Actor loss (policy gradient with advantage)
        action_probs = self.actor(states)
        log_probs = torch.log(action_probs.gather(1, actions.unsqueeze(1)).squeeze())
        actor_loss = -(log_probs * advantages.detach()).mean()
        
        # Add entropy bonus (encourage exploration)
        dist = torch.distributions.Categorical(action_probs)
        entropy = dist.entropy().mean()
        actor_loss -= self.entropy_coef * entropy
        
        # Critic loss (MSE between value and TD target)
        critic_loss = nn.MSELoss()(values, td_targets.detach())
        
        # Update networks
        self.actor_optimizer.zero_grad()
        actor_loss.backward()
        torch.nn.utils.clip_grad_norm_(self.actor.parameters(), 1.0)
        self.actor_optimizer.step()
        
        self.critic_optimizer.zero_grad()
        critic_loss.backward()
        torch.nn.utils.clip_grad_norm_(self.critic.parameters(), 1.0)
        self.critic_optimizer.step()
        
        return {
            'actor_loss': actor_loss.item(),
            'critic_loss': critic_loss.item(),
            'entropy': entropy.item(),
            'advantage_mean': advantages.mean().item()
        }
\end{lstlisting}

\section{Data Collection Pipeline}

\subsection{Real PubMed Data Collection}

Training data collected from live PubMed API.

\textbf{Collection Script:}
\begin{lstlisting}[style=bashstyle]
# rl_optimization/data_collection/collect_real_pubmed_data.py
python collect_real_pubmed_data.py \
    --max-results 500 \
    --queries "adverse events" "drug safety" "pharmacovigilance" "case report" \
    --output rl_optimization/real_pubmed_data \
    --enable-classification
\end{lstlisting}

\textbf{Data Structure:}
\begin{lstlisting}[style=pythonstyle]
# rl_optimization/real_pubmed_data/collected_data.json
{
  "documents": [
    {
      "pmid": "12345678",
      "title": "Adverse Events Following Aspirin Use",
      "abstract": "...",
      "content": "...",
      "mesh_terms": ["Aspirin", "Adverse Drug Reaction"],
      "publication_types": ["Case Reports"],
      "classification": "ICSR",  # Ground truth from ensemble
      "classification_confidence": 0.94,
      "ensemble_metadata": {
        "rule_based": {"class": "ICSR", "confidence": 0.89},
        "bert": {"class": "Safety_Information", "confidence": 0.72},
        "llm": {"class": "ICSR", "confidence": 0.96}
      },
      "processing_time_ms": 2340,
      "cost": 0.0001
    }
  ],
  "statistics": {
    "total_documents": 500,
    "icsr_count": 120,
    "safety_info_count": 180,
    "non_icsr_count": 150,
    "rejected_count": 50
  }
}
\end{lstlisting}

\subsection{Feature Extraction}

Extracting 47-dimensional state vectors from collected data.

\begin{lstlisting}[style=bashstyle]
# Extract features for RL training
python rl_optimization/data_collection/feature_extractor.py \
    --input rl_optimization/real_pubmed_data/collected_data.json \
    --output rl_optimization/real_pubmed_data/feature_vectors.npy

# Output: 500 x 47 numpy array
\end{lstlisting}

\section{Training Process}

\subsection{Training Configuration}

\begin{lstlisting}[style=pythonstyle]
# rl_optimization/config/training_config.yaml
training:
  episodes: 1000
  max_steps_per_episode: 50
  batch_size: 32
  learning_rate: 0.001
  gamma: 0.99
  entropy_coef: 0.01

data:
  train_split: 0.8  # 400 documents
  val_split: 0.2    # 100 documents
  shuffle: true

optimization:
  optimizer: adam
  gradient_clip: 1.0
  early_stopping:
    patience: 50
    min_delta: 0.001

checkpointing:
  save_interval: 100  # Save every 100 episodes
  keep_best: true
  checkpoint_dir: rl_optimization/checkpoints

logging:
  log_interval: 10
  tensorboard: true
  mlflow: true
\end{lstlisting}

\subsection{Training Script}

\begin{commandbox}
\begin{lstlisting}[style=bashstyle]
# Start RL training
python rl_optimization/rl_training/training_engine.py \
    --config rl_optimization/config/training_config.yaml \
    --data-path rl_optimization/real_pubmed_data \
    --output rl_optimization/training_results

# Training output:
Episode 1/1000: Reward=0.12, Loss=0.45, Accuracy=75%
Episode 10/1000: Reward=0.24, Loss=0.38, Accuracy=78%
Episode 50/1000: Reward=0.41, Loss=0.25, Accuracy=82%
...
Episode 987/1000: Reward=0.84, Loss=0.08, Accuracy=92.8%  # PEAK!
Episode 1000/1000: Reward=0.82, Loss=0.09, Accuracy=92.1%

Training Complete!
Best Episode: 987 (Reward: 0.84, Accuracy: 92.8%)
Total Training Time: 4.2 hours
\end{lstlisting}
\end{commandbox}

\subsection{Training Results}

\begin{figure}[h]
\centering
\fbox{\textit{[SCREENSHOT PLACEHOLDER: RL Training Curves - Reward, Loss, Accuracy over Episodes]}}
\caption{RL Training Progress Over 1000 Episodes}
\label{fig:rl-training-curves}
\end{figure}

\textbf{Performance Improvements:}
\begin{table}[h]
\centering
\caption{RL Training Results - Before vs. After}
\begin{tabularx}{\textwidth}{lXXXX}
\toprule
\textbf{Metric} & \textbf{Baseline} & \textbf{Episode 1} & \textbf{Episode 987 (Peak)} & \textbf{Improvement} \\
\midrule
Accuracy & 87.2\% & 75.0\% & 92.8\% & +5.6\% \\
Cost per 1K docs & \$8.40 & \$12.50 & \$2.95 & -65\% \\
Processing Time & 3.2s & 4.1s & 1.8s & -44\% \\
Confidence Variance & 0.28 & 0.42 & 0.11 & -61\% \\
Total Reward & 0.34 & 0.12 & 0.84 & +147\% \\
\bottomrule
\end{tabularx}
\end{table}

\textbf{Learned Model Weights:}
\begin{itemize}
    \item Rule-based: 0.30 → \textbf{0.48} (+60\%)
    \item BERT: 0.48 → \textbf{0.30} (-38\%)
    \item LLM: 0.22 → \textbf{0.27} (+23\%)
\end{itemize}

\textbf{Key Insights:}
\begin{enumerate}
    \item RL learned to \textbf{trust rule-based classifier more} (highest confidence when rules are certain)
    \item RL learned to \textbf{penalize BERT} (least reliable in PV domain)
    \item RL learned to \textbf{selectively invoke LLM} (only for complex cases)
    \item This strategy maximizes accuracy while minimizing cost
\end{enumerate}

\section{Model Deployment}

\subsection{Model Inference}

\textbf{Load Trained Model:}
\begin{lstlisting}[style=pythonstyle]
# services/rl-intelligence-service/main.py
import torch

class RLInferenceService:
    def __init__(self):
        self.agent = self.load_trained_agent()
    
    def load_trained_agent(self):
        """Load trained RL model"""
        checkpoint_path = "rl_optimization/checkpoints/best_model.pt"
        checkpoint = torch.load(checkpoint_path)
        
        agent = A2CAgent(state_dim=47, action_dim=12)
        agent.actor.load_state_dict(checkpoint['actor_state_dict'])
        agent.critic.load_state_dict(checkpoint['critic_state_dict'])
        agent.actor.eval()  # Set to evaluation mode
        agent.critic.eval()
        
        return agent
    
    def get_optimal_configuration(self, document: dict) -> dict:
        """Get RL-optimized classification configuration"""
        # Extract state from document
        state = self.feature_extractor.extract_state(document)
        
        # Get optimal action from trained policy
        action, _, _ = self.agent.select_action(state, training=False)
        
        # Apply action to base configuration
        config = self.agent.apply_action(action)
        
        return config
\end{lstlisting}

\subsection{Integration with Classification Service}

\textbf{RL-Enhanced Classification:}
\begin{lstlisting}[style=pythonstyle]
# services/pv-classification-service/main.py
async def classify_with_rl_optimization(document: dict):
    """Classify document using RL-optimized configuration"""
    
    # Get RL-optimized configuration
    rl_config = await rl_service.get_optimal_configuration(document)
    
    # Use optimized weights for ensemble
    ensemble = EnsembleClassifier(
        rule_weight=rl_config['rule_weight'],
        bert_weight=rl_config['bert_weight'],
        llm_weight=rl_config['llm_weight'],
        rule_threshold=rl_config['rule_threshold'],
        bert_threshold=rl_config['bert_threshold'],
        cascade_enabled=rl_config['cascade_enabled']
    )
    
    # Classify
    result = ensemble.classify(document)
    
    return result
\end{lstlisting}

\section{Monitoring \& Evaluation}

\subsection{RL Intelligence Dashboard}

Real-time monitoring of RL performance.

\textbf{Key Metrics Tracked:}
\begin{itemize}
    \item Current accuracy (rolling 100 documents)
    \item Cost savings vs. baseline
    \item Processing speed improvement
    \item Configuration adjustments per hour
    \item Reward trend (smoothed)
    \item Action distribution
\end{itemize}

\begin{figure}[h]
\centering
\fbox{\textit{[SCREENSHOT PLACEHOLDER: RL Intelligence Dashboard - Real-time Metrics]}}
\caption{RL Intelligence Dashboard - Live Performance}
\label{fig:rl-dashboard-live}
\end{figure}

\subsection{Continuous Learning (Future)}

\textbf{Online Learning Pipeline (Roadmap):}
\begin{enumerate}
    \item Collect user feedback on classifications
    \item Store (state, action, reward) tuples in replay buffer
    \item Periodically retrain agent (weekly batch updates)
    \item A/B test new model vs. current model
    \item Deploy if performance improves by 2\%+
\end{enumerate}

\end{document}

